\chapter{BDD-тестирование} \label{ch1}

В данной главе описывается план BDD-тестирования консольного приложения поиска палиндромных подстрок. Рассматриваются причины эволюции TDD в BDD, синтаксис Gherkin, выбор фреймворка и набор разработанных пользовательских историй.

\section{Основы BDD и Gherkin} \label{ch1:sec1}

В статье Дэна Норта BDD предлагается как ответ на практические трудности TDD: разработчики часто спорят о том, что именно считать тестом, как называть тестовые методы и почему технические тесты плохо объясняют ценность функциональности для пользователя~\cite{north-bdd}. BDD сохраняет идею ранней формализации ожидаемого результата, но описывает проверку как поведение системы в бизнес-контексте.

Основная единица BDD --- сценарий, записанный на языке Gherkin. Такой сценарий читается как спецификация и одновременно может выполняться автоматизированным фреймворком. Ключевые конструкции Gherkin приведены в \taref{tab:gherkin}.

\noindent
\begingroup
\centering
\small
\begin{longtable}[c]{|p{3.2cm}|p{11.8cm}|}
    \caption{Основные конструкции Gherkin}
    \label{tab:gherkin}
    \\
    \hline
    Конструкция & Назначение \\ \hline
    \endfirsthead
    \hline
    Конструкция & Назначение \\ \hline
    \endhead
    \hline
    \endfoot
    Feature / Функционал & Описывает группу сценариев и пользовательскую историю: кто, что хочет сделать и какую получает выгоду \\ \hline
    Rule / Правило & Фиксирует бизнес-правило, объединяющее несколько сценариев внутри одной функциональности \\ \hline
    Background / Предыстория & Задаёт общие предусловия, выполняемые перед сценариями \\ \hline
    Scenario / Сценарий & Описывает один конкретный пример поведения системы \\ \hline
    Scenario Outline / Структура сценария & Позволяет параметризовать сценарий таблицей Examples \\ \hline
    Given / Допустим & Описывает исходное состояние системы \\ \hline
    When / Когда / Если & Описывает действие пользователя или внешнее событие \\ \hline
    Then / Тогда / То & Описывает ожидаемый наблюдаемый результат \\ \hline
    And, But / И, Но & Дополняют предыдущий шаг без изменения его логического типа \\ \hline
\end{longtable}
\normalsize
\endgroup

В работе используются как декларативные сценарии, описывающие желаемый результат без деталей интерфейса, так и императивный сценарий, фиксирующий последовательность действий пользователя в меню. Такое сочетание соответствует требованию задания и позволяет проверить поведение приложения на разных уровнях абстракции.

\section{Выбор BDD-фреймворка} \label{ch1:sec2}

Проект реализован на JavaScript в среде Node.js, поэтому в качестве BDD-фреймворка выбран Cucumber.js~\cite{cucumber-js}. Причины выбора:
\begin{itemize}
    \item поддерживает Gherkin и локализацию ключевых слов на русский язык;
    \item позволяет хранить сценарии в отдельных \texttt{.feature}-файлах;
    \item связывает шаги сценариев с JavaScript-функциями через Step Definitions;
    \item интегрируется с существующим проектом без изменения основной архитектуры;
    \item формирует HTML-отчёт о выполнении BDD-сценариев.
\end{itemize}

Команда запуска BDD-тестов добавлена в \texttt{package.json}:
\begin{center}
\texttt{npm run bdd}
\end{center}

Она запускает \texttt{cucumber-js}, подключает \texttt{features/step\_definitions} и \texttt{features/support}, а также сохраняет HTML-отчёт в каталог \texttt{reports/cucumber}.

\section{Описание тестовых случаев и BDD-сценариев} \label{ch1:sec3}

Для приложения разработаны три Gherkin-истории, каждая размещена в отдельном \texttt{.feature}-файле:
\begin{enumerate}
    \item \texttt{palindrome\_analysis.feature} --- анализ палиндромных подстрок;
    \item \texttt{json\_io.feature} --- загрузка входной строки и сохранение результата;
    \item \texttt{performance\_comparison.feature} --- нагрузочное сравнение Манакера и расширения от центра.
\end{enumerate}

Сводная таблица покрытия приведена в \taref{tab:bdd-coverage}.

\noindent
\begingroup
\centering
\small
\begin{longtable}[c]{|p{4.2cm}|p{4.3cm}|p{6.3cm}|}
    \caption{Сводная таблица BDD-покрытия}
    \label{tab:bdd-coverage}
    \\
    \hline
    Файл & Проверяемый модуль & Основные элементы задания \\ \hline
    \endfirsthead
    \hline
    Файл & Проверяемый модуль & Основные элементы задания \\ \hline
    \endhead
    \hline
    \endfoot
    \texttt{palindrome\_analysis.feature} & \texttt{palindromeService}, алгоритмы & Feature, Rule, Background, Scenario Outline, таблица Examples, декларативный и императивный стиль \\ \hline
    \texttt{json\_io.feature} & \texttt{jsonReader}, \texttt{resultWriter} & пользовательская история, Given/When/Then/And/But, негативный сценарий \\ \hline
    \texttt{performance\_comparison.feature} & генераторы данных, алгоритмы & сравнение двух алгоритмов, элементы нагрузочного тестирования, входная таблица \\ \hline
\end{longtable}
\normalsize
\endgroup

\section{Примеры сценариев} \label{ch1:sec4}

Декларативный сценарий из \texttt{palindrome\_analysis.feature} проверяет итоговое поведение без описания внутренних вызовов:

\begin{lstlisting}[breaklines=true]
Структура сценария: Декларативная проверка результата для строк разных типов
  Если пользователь анализирует строку "<input>" алгоритмом Манакера
  То самый длинный палиндром равен "<longest>"
  И количество палиндромных подстрок равно <count>
  И результат совпадает с эталонным алгоритмом расширения от центра
\end{lstlisting}

Императивный сценарий в том же файле описывает последовательность действий меню:

\begin{lstlisting}[breaklines=true]
Сценарий: Императивная проверка запуска через последовательность действий меню
  Допустим пользователь ввел строку "abbaabba"
  Когда пользователь выбирает пункт "Запустить алгоритм Манакера"
  И пользователь выбирает пункт "Показать подробный отчет"
  То отчет содержит строку "abbaabba"
\end{lstlisting}

В сценарии производительности используются два алгоритма решения наукоемкой задачи: алгоритм Манакера и расширение от центра. Для каждого набора данных проверяется совпадение результата, а время Манакера ограничивается заданным порогом.

\section{Сопоставление BDD и модульных тестов} \label{ch1:sec5}

В проекте сохраняются существующие модульные тесты Jest, так как они проверяют детали вычисления радиусов, граничные условия и ошибки ввода-вывода. BDD-сценарии решают другую задачу: они описывают поведение приложения на языке требований. Например, unit-тест может проверять точный массив \texttt{d\_odd}, а BDD-сценарий проверяет, что пользователь получает правильный самый длинный палиндром и подтверждение совпадения двух алгоритмов.

Таким образом, BDD не заменяет модульные тесты, а надстраивается над ними как уровень исполняемой спецификации.

\section{Выводы} \label{ch1:conclusion}

В результате для приложения подготовлен план BDD-тестирования и три Gherkin-истории. В них присутствуют пользовательские истории, \texttt{Rule}, \texttt{Background}, \texttt{Scenario Outline}, таблицы входных данных, шаги \texttt{Given}, \texttt{When}, \texttt{Then}, \texttt{And}, \texttt{But}, а также сценарий нагрузочного сравнения двух алгоритмов. Выбранный фреймворк Cucumber.js соответствует языку реализации проекта и позволяет выполнять спецификации как автоматизированные тесты.
